\documentclass{article}
\usepackage{graphicx}
\usepackage[T1]{fontenc}
\usepackage[utf8]{inputenc}
\usepackage{lmodern}
\usepackage{geometry}
\geometry{margin=1in}
\usepackage{microtype}
\usepackage{amsmath,amssymb,amsfonts}
\usepackage{mathtools}
\usepackage{booktabs}
\usepackage{siunitx}
\usepackage{sectsty}
\usepackage{enumitem}
\sisetup{detect-all=true,round-mode=places,round-precision=3}

\numberwithin{equation}{subsection}

\title{Value of Forecasting Accuracy --- Day-Ahead Auction Only \\
\large Methodology v1.1: Round-Trip Efficiency, Opportunity-Cost Bidding,\\
and Volume-Balanced Paired Rejection}
\author{Erwin Spolders, FlexUp GmbH}
\date{April 2026}

\begin{document}

\maketitle

\section{Introduction}

\subsection*{Overview}
This document presents the v1.1 mathematical formulation of the CastValue framework: a comparison of the energy arbitrage margins of a Battery Energy Storage System (BESS) in the Day-Ahead Auction (DAA) power market in Germany, under different price forecasts.

v1.1 extends the v1.0 formulation (see \texttt{DAA\_DE\_LU\_Formulation.pdf}) with the physical and market frictions a real trader faces:

\begin{itemize}[itemsep=2pt, topsep=2pt]
    \item \textbf{Round-Trip Efficiency (RTE)} of the BESS in the LP energy balance --- one MWh delivered to the grid requires more than one MWh purchased from the grid.
    \item \textbf{Opportunity-cost bid derivation} --- limit prices for charge bids and discharge asks are derived from the forecast and adjusted for RTE.
    \item \textbf{Out-of-the-money (OTM) rejection} --- bids/asks not in the money against the actual clearing price do not fill.
    \item \textbf{Volume-balanced paired rejection} --- when OTM fills leave cleared charge and discharge volumes unmatched, additional MTUs are dropped on the oversubscribed side until volumes balance and end-of-day SOC equals start-of-day SOC.
    \item \textbf{Rolling SOC} across days --- the day's initial SOC equals the previous day's terminal SOC (the LP is no longer forced to start and end at zero).
\end{itemize}

The Perfect Foresight benchmark in v1.1 is defined as the LP optimum given RTE, with bids placed at the clearing price (100\% fill). It is a ``perfect scheduling'' benchmark, not a frictionless one --- the $(1-\eta)$ round-trip loss is fully paid in the realised margin.

The v1.1 dataset covers EPEX SPOT DE-LU at native quarter-hourly resolution from 18 October 2025 to 22 February 2026 (126 days). Because the dataset is entirely post the 1~October 2025 market structure change, the hourly coupling constraint of v1.0 is dropped.

This document accompanies the Python implementation of the v1.1 optimization available on our GitHub.


\section*{Symbols}
Notation extends v1.0. Symbols introduced in v1.1 are flagged in the description column.

\subsection*{Sets and indices}
\begin{tabular}{@{}ll@{}}
\toprule
Symbol & Meaning \\
\midrule
\(D=\{1,\dots,126\}\) & day index (18~Oct~2025 -- 22~Feb~2026) \\
\(Q=\{1,\dots,96\}\) & quarter-hour index within the day \\
\(Q^{+}=\{1,\dots,97\}\) & quarter-hour index for SOC including terminal state \\
\(\mathcal{K}=\{\mathrm{P},1,\dots,k,\mathrm{PF}\}\) & price forecast index (Persistence, forecasts \(1\)\dots\(k\), Perfect Foresight) \\
\(\hat{\lambda}^{(\mathrm{P})}\) & Persistence forecast, where \(\hat{\lambda}^{(\mathrm{P})}_{d,q}:=\lambda_{d-1,q}\) \\
\(\hat{\lambda}^{(\mathrm{PF})}\) & Perfect Foresight forecast, where \(\hat{\lambda}^{(\mathrm{PF})}_{d,q}:=\lambda_{d,q}\) \\
\bottomrule
\end{tabular}

\subsection*{Parameters}
\begingroup
\renewcommand{\arraystretch}{1.25}
\begin{tabular}{@{}lll@{}}
\toprule
Symbol & Meaning & Notes \\
\midrule
\(\texttt{cap}^{\text{energy}}\) & energy capacity [MWh] & \\
\(\texttt{cap}^{\text{power}}\) & charge/discharge power limit [MW] & \\
\(n_{\text{cycles}}\) & allowed daily equivalent full cycles [--] & \\
\(L_d =\texttt{cap}^{\text{energy}}\cdot n_{\text{cycles}}\) & daily energy-throughput limit [MWh] & \\
\(\hat{\lambda}^{\text{DAA, k}}_{d,q}\) & DAA forecast price from forecast \(k\) [EUR/MWh] & \\
\(\lambda^{\text{DAA}}_{d,q}\) & DAA actual clearing price [EUR/MWh] & \\
\(\Delta=\tfrac{1}{4}\) & quarter-hour duration [h] & \\
\(\eta\) & round-trip efficiency [--] (\(\eta=0.85\)) & \textbf{new in v1.1} \\
\bottomrule
\end{tabular}
\endgroup

\subsection*{LP decision variables}
\begingroup
\renewcommand{\arraystretch}{1.5}
\begin{tabular}{@{}ll@{}}
\toprule
Symbol & Meaning \\
\midrule
\(\mathrm{CHA}^{\text{DAA}}_{q,k},\ \mathrm{DIS}^{\text{DAA}}_{q,k} \in [0,1]\) & normalized planned charge/discharge in DAA at \(q,k\) \\
\(\mathrm{SOC}_{q,k}\) & planned state-of-charge at \(q\in Q^+\) [MWh] \\
\bottomrule
\end{tabular}
\endgroup

\subsection*{Derived quantities (post-LP, v1.1)}
\begingroup
\renewcommand{\arraystretch}{1.5}
\begin{tabular}{@{}ll@{}}
\toprule
Symbol & Meaning \\
\midrule
\(\bar{\lambda}^{\text{cha}}_{d,k},\ \bar{\lambda}^{\text{dis}}_{d,k}\) & volume-weighted average forecast charge / discharge price [EUR/MWh] \\
\(B^{\text{cha}}_{q,k},\ B^{\text{dis}}_{q,k}\) & charge bid / discharge ask price [EUR/MWh] \\
\(F^{\text{cha}}_{q,k},\ F^{\text{dis}}_{q,k} \in \{0,1\}\) & OTM fill indicators \\
\(\widetilde{\mathrm{CHA}}_{q,k},\ \widetilde{\mathrm{DIS}}_{q,k}\) & post-OTM accepted volumes [normalized] \\
\(\mathrm{CHA}^{*}_{q,k},\ \mathrm{DIS}^{*}_{q,k}\) & realised (delivered) volumes after paired rejection [normalized] \\
\(\mathrm{SOC}^{*}_{q,d,k}\) & realised state-of-charge trajectory [MWh] \\
\bottomrule
\end{tabular}
\endgroup

\subsection*{Outputs}
\begin{tabular}{@{}ll@{}}
\toprule
Symbol & Meaning \\
\midrule
\(V_k\) & Value of forecast \(k\) compared with persistence forecast \(\mathrm{P}\) \\
\(C_k\) & Cost of the error of forecast \(k\) compared with perfect foresight \(\mathrm{PF}\) \\
\(AR_{\text{total}, k}\) & Total trading margin from DAA across all days \(d \in D\), given forecast \(k\) \\
\(DR_{d, k}\) & Daily trading margin from DAA for day \(d\), given forecast \(k\) \\
\bottomrule
\end{tabular}


\section{Calculating Outputs}

\subsection{Forecast Value}
The value (per MW traded BESS) of a price forecast equals the difference in the trading margin generated by trading on that forecast vs.\ the persistence forecast:
\begin{equation}
V_k = AR_{\text{total},k} - AR_{\text{total},P}
\end{equation}

\subsection{Cost of Forecast Error}
The cost of forecast error equals the gap between Perfect Foresight and forecast \(k\):
\begin{equation}
C_k = AR_{\text{total}, PF} - AR_{\text{total},k}
\end{equation}

\subsection{Total and Daily Trading Margin}
\begin{equation}
AR_{\text{total}, k} = \sum_{d \in D} DR_{d, k}
\end{equation}

The daily trading margin uses \emph{realised} delivered volumes, evaluated at actual prices:
\begin{equation}
DR_{d, k} = \sum_{q\in Q} \lambda^{\text{DAA}}_{d,q} \cdot \texttt{cap}^{\text{power}} \cdot \Delta \cdot
\bigl(\mathrm{DIS}^{*}_{q,d,k} - \mathrm{CHA}^{*}_{q,d,k}\bigr)
\label{eq:daily-margin}
\end{equation}

The factor \(\texttt{cap}^{\text{power}}\cdot\Delta\) converts normalized values to grid-side MWh per quarter-hour. Note the trader pays for the full grid-side charge volume, including the \((1-\eta)\) portion lost to RTE.


\section{Optimizing Daily Trading Decisions (LP)}

\subsection{Day-Ahead Auction (DAA) --- v1.1}

\subsection*{Objective}
Maximise expected revenue at forecast prices, given decisions \(\mathrm{CHA}^{\mathrm{DAA}}_{q,k}, \mathrm{DIS}^{\mathrm{DAA}}_{q,k}\):
\begin{equation}
\operatorname*{max}_{\mathrm{CHA}^{\mathrm{DAA}}_{q,k},\,\mathrm{DIS}^{\mathrm{DAA}}_{q,k}}
\sum_{q\in Q} \hat{\lambda}^{\mathrm{DAA}}_{q,k} \cdot \texttt{cap}^{\text{power}} \cdot \Delta \cdot
\bigl(\mathrm{DIS}^{\mathrm{DAA}}_{q,k}-\mathrm{CHA}^{\mathrm{DAA}}_{q,k}\bigr)
\end{equation}

\subsection*{Constraints}
SOC bounds:
\begin{equation}
0 \leq \mathrm{SOC}_q \leq \texttt{cap}^{\text{energy}} \quad \forall q \in Q^+
\end{equation}

\textbf{Rolling initial SOC} (v1.1 replaces the v1.0 condition \(\mathrm{SOC}_1 = 0\)):
\begin{equation}
\mathrm{SOC}_{1,d,k} = \mathrm{SOC}^{*}_{97,\,d-1,\,k} \qquad \text{with } \mathrm{SOC}^{*}_{97,\,0,\,k} := 0
\end{equation}

\textbf{Terminal SOC} non-negative but free (v1.1 replaces the v1.0 condition \(\mathrm{SOC}_{97} = 0\)):
\begin{equation}
\mathrm{SOC}_{97} \geq 0
\end{equation}

\textbf{RTE-aware energy balance} (v1.1 generalises the v1.0 step constraint):
\begin{equation}
\mathrm{SOC}_{q+1} = \mathrm{SOC}_{q}
+ \texttt{cap}^{\text{power}} \cdot \Delta \cdot
\bigl(\eta \cdot \mathrm{CHA}^{\mathrm{DAA}}_q - \mathrm{DIS}^{\mathrm{DAA}}_q \bigr)
\quad \forall q \in Q
\label{eq:soc-step}
\end{equation}
i.e.\ each MWh purchased from the grid stores only \(\eta\) MWh in the battery; each MWh delivered to the grid draws 1~MWh from the battery.

\textbf{Volume caps} (battery-side cycles): the discharge cap remains \(L_d\), and the charge cap is relaxed by \(1/\eta\) so the LP can deliver a full discharge cycle despite RTE losses:
\begin{align}
\texttt{cap}^{\text{power}} \cdot \Delta \cdot \sum_{q \in Q} \mathrm{CHA}^{\mathrm{DAA}}_q &\leq \frac{L_d}{\eta} \\
\texttt{cap}^{\text{power}} \cdot \Delta \cdot \sum_{q \in Q} \mathrm{DIS}^{\mathrm{DAA}}_q &\leq L_d
\end{align}

The hourly coupling constraint of v1.0 is dropped (the v1.1 dataset is entirely 15-minute).


\section{Bid Derivation}

The trader posts limit prices for each LP-active MTU using a volume-weighted opportunity-cost rule.

\subsection{Volume-weighted average forecast prices}

Define the LP-planned grid-side energies:
\begin{equation}
E^{\text{cha}}_{q,k} := \texttt{cap}^{\text{power}} \cdot \Delta \cdot \mathrm{CHA}^{\mathrm{DAA}}_{q,k},
\qquad
E^{\text{dis}}_{q,k} := \texttt{cap}^{\text{power}} \cdot \Delta \cdot \mathrm{DIS}^{\mathrm{DAA}}_{q,k}.
\end{equation}

The volume-weighted average forecast charge and discharge prices for day \(d\), forecast \(k\) are:
\begin{equation}
\bar{\lambda}^{\text{cha}}_{d,k}
= \frac{\sum_{q\in Q}\hat{\lambda}^{\text{DAA}}_{q,k}\,E^{\text{cha}}_{q,k}}
       {\sum_{q\in Q}E^{\text{cha}}_{q,k}},
\qquad
\bar{\lambda}^{\text{dis}}_{d,k}
= \frac{\sum_{q\in Q}\hat{\lambda}^{\text{DAA}}_{q,k}\,E^{\text{dis}}_{q,k}}
       {\sum_{q\in Q}E^{\text{dis}}_{q,k}}.
\end{equation}

\subsection{Bid prices}

For LP-active MTUs (\(\mathrm{CHA}^{\mathrm{DAA}}_{q,k} > 0\) or \(\mathrm{DIS}^{\mathrm{DAA}}_{q,k} > 0\)):
\begin{align}
B^{\text{cha}}_{q,k} &= \eta \cdot \bar{\lambda}^{\text{dis}}_{d,k}
\qquad\text{if } \mathrm{CHA}^{\mathrm{DAA}}_{q,k} > 0 \\
B^{\text{dis}}_{q,k} &= \frac{\bar{\lambda}^{\text{cha}}_{d,k}}{\eta}
\qquad\text{if } \mathrm{DIS}^{\mathrm{DAA}}_{q,k} > 0
\end{align}

\textbf{Interpretation.} The trader pays no more for charging than the expected revenue from later discharging \emph{after} RTE loss; and asks no less for discharging than the cost of the energy charged \emph{gross of} RTE loss. These are the indifference prices implied by the day's forecast schedule.


\section{Out-of-the-Money (OTM) Fill Rule}

A bid fills only if it is in the money against the actual clearing price:
\begin{align}
F^{\text{cha}}_{q,k} &= \mathbf{1}\bigl[\,\mathrm{CHA}^{\mathrm{DAA}}_{q,k} > 0 \;\wedge\; B^{\text{cha}}_{q,k} \geq \lambda^{\text{DAA}}_{q,d}\,\bigr] \\
F^{\text{dis}}_{q,k} &= \mathbf{1}\bigl[\,\mathrm{DIS}^{\mathrm{DAA}}_{q,k} > 0 \;\wedge\; B^{\text{dis}}_{q,k} \leq \lambda^{\text{DAA}}_{q,d}\,\bigr]
\end{align}

Post-OTM accepted volumes:
\begin{equation}
\widetilde{\mathrm{CHA}}_{q,k} = F^{\text{cha}}_{q,k} \cdot \mathrm{CHA}^{\mathrm{DAA}}_{q,k},
\qquad
\widetilde{\mathrm{DIS}}_{q,k} = F^{\text{dis}}_{q,k} \cdot \mathrm{DIS}^{\mathrm{DAA}}_{q,k}.
\end{equation}


\section{Volume-Balanced Paired Rejection}

When forecast errors lead to asymmetric OTM fills, cleared charge and discharge volumes diverge. To enforce energy balance and avoid open positions at end-of-day, the worst-margin MTUs on the oversubscribed side are progressively rejected until volumes match.

Define the cleared battery-side volumes after OTM:
\begin{align}
V^{\text{cha-stored}}_{d,k} &= \texttt{cap}^{\text{power}} \cdot \Delta \cdot \eta \cdot \sum_{q\in Q} \widetilde{\mathrm{CHA}}_{q,k} \\
V^{\text{dis-drawn}}_{d,k} &= \texttt{cap}^{\text{power}} \cdot \Delta \cdot \sum_{q\in Q} \widetilde{\mathrm{DIS}}_{q,k}
\end{align}

\textbf{Procedure.} Initialise \(\mathrm{CHA}^{*}_{q,k} \leftarrow \widetilde{\mathrm{CHA}}_{q,k}\) and \(\mathrm{DIS}^{*}_{q,k} \leftarrow \widetilde{\mathrm{DIS}}_{q,k}\).

\begin{enumerate}[itemsep=2pt, topsep=4pt]
    \item If \(V^{\text{cha-stored}}_{d,k} > V^{\text{dis-drawn}}_{d,k}\): identify cleared charge MTUs and sort by descending actual price \(\lambda^{\mathrm{DAA}}_{q,d}\) (most-expensive purchase first = worst margin first). Sequentially reduce \(\mathrm{CHA}^{*}_{q,k}\) until battery-side volumes match.
    \item If \(V^{\text{dis-drawn}}_{d,k} > V^{\text{cha-stored}}_{d,k}\): identify cleared discharge MTUs and sort by ascending actual price (cheapest sale first = worst margin first). Sequentially reduce \(\mathrm{DIS}^{*}_{q,k}\) until battery-side volumes match.
    \item Otherwise (\(V^{\text{cha-stored}}_{d,k} = V^{\text{dis-drawn}}_{d,k}\)): no adjustment needed.
\end{enumerate}

The realised SOC trajectory is then computed analogously to \eqref{eq:soc-step}, using the realised volumes:
\begin{equation}
\mathrm{SOC}^{*}_{q+1,d,k} = \mathrm{SOC}^{*}_{q,d,k}
+ \texttt{cap}^{\text{power}} \cdot \Delta \cdot
\bigl(\eta \cdot \mathrm{CHA}^{*}_{q,k} - \mathrm{DIS}^{*}_{q,k}\bigr)
\quad \forall q \in Q
\end{equation}
with \(\mathrm{SOC}^{*}_{1,d,k} = \mathrm{SOC}^{*}_{97,d-1,k}\). By construction \(\mathrm{SOC}^{*}_{97,d,k} = \mathrm{SOC}^{*}_{1,d,k}\) --- daily energy balances, and there is no carryover or imbalance position.


\section{Perfect Foresight Benchmark}

Under Perfect Foresight (\(\hat{\lambda}^{\mathrm{DAA}}_{q,PF} = \lambda^{\mathrm{DAA}}_{q,d}\)), the trader knows the clearing price exactly and posts limit prices at clearing:
\begin{equation}
B^{\text{cha}}_{q,PF} = B^{\text{dis}}_{q,PF} = \lambda^{\text{DAA}}_{q,d}.
\end{equation}
All LP-scheduled MTUs are then in the money (\(F^{\text{cha}}_{q,PF} = F^{\text{dis}}_{q,PF} = 1\) wherever the LP is active), so:
\begin{equation}
\mathrm{CHA}^{*}_{q,PF} = \mathrm{CHA}^{\mathrm{DAA}}_{q,PF},
\qquad
\mathrm{DIS}^{*}_{q,PF} = \mathrm{DIS}^{\mathrm{DAA}}_{q,PF}.
\end{equation}
Paired rejection is a no-op. The Perfect Foresight margin equals the LP optimum subject to RTE --- not a frictionless ceiling. The \((1-\eta)\) round-trip loss is fully reflected in \eqref{eq:daily-margin}.


\section{Summary of Differences from v1.0}

\begin{tabular}{@{}p{0.36\linewidth}p{0.28\linewidth}p{0.32\linewidth}@{}}
\toprule
Aspect & v1.0 & v1.1 \\
\midrule
RTE in LP & none (1:1 efficiency) & \(\eta=0.85\) on charge side \\
Charge-side volume cap & \(L_d\) & \(L_d / \eta\) \\
Initial SOC & \(\mathrm{SOC}_1 = 0\) & rolling: \(\mathrm{SOC}_{1,d} = \mathrm{SOC}^{*}_{97,d-1}\) \\
Terminal SOC & \(\mathrm{SOC}_{97} = 0\) & \(\mathrm{SOC}_{97} \geq 0\) (free) \\
Hourly coupling & active before 1~Oct~2025 & dropped (15-min only) \\
Bid pricing & implicit (zero-cap, infinite-floor) & opportunity-cost, RTE-adjusted \\
Bid acceptance & 100\% (guaranteed) & OTM check vs.\ actual \\
Volume reconciliation & none (LP guarantees by \(\mathrm{SOC}_{97}=0\)) & paired rejection \\
Perfect Foresight & frictionless LP optimum & RTE-aware LP optimum, bids at clearing \\
\bottomrule
\end{tabular}

\end{document}
